\documentclass[10pt]{article}

% ---------- Document Identity (update these only) ----------
\newcommand{\DocStage}{}
\newcommand{\DocTitle}{Your Road to Tier One \textbf{SOF}}
\newcommand{\ProgramName}{182-Week Civilian-to-Selection Readiness Pipeline}
\newcommand{\ProgramSubtitle}{}

% ---------- Page ----------
\usepackage[legalpaper,left=0.5in,right=0.5in,top=0.6in,bottom=0.45in]{geometry}

% ---------- Typography ----------
\usepackage[T1]{fontenc}
\usepackage[utf8]{inputenc}
\usepackage{libertinus}
\usepackage{microtype}
\usepackage{parskip}
\usepackage{ragged2e}
\usepackage{textcomp}
\AtBeginDocument{\color{black}}
\AtBeginDocument{\pagecolor{white}}
\AtBeginDocument{\justifying}

% ---------- Landscape pages ----------
\usepackage{pdflscape}

% ---------- Color ----------
\usepackage[table]{xcolor}
\definecolor{StageBlue}{HTML}{1F4E79}
\definecolor{StageBlueDark}{HTML}{163A5A}
\definecolor{LightGray}{HTML}{F2F4F7}
\definecolor{MidGray}{HTML}{D9DEE5}
\definecolor{RuleGray}{HTML}{C7CED8}
\definecolor{HeaderInk}{HTML}{000000}
\definecolor{Ink}{HTML}{000000}

% ---------- Utility ----------
\usepackage{enumitem}
\sloppy
\setlength{\emergencystretch}{2em}

% ---------- Boxes / UI ----------
\usepackage[most]{tcolorbox}

\tcbset{
  charterTitle/.style={
    enhanced,
    colback=StageBlue!4,
    colframe=StageBlue,
    boxrule=1.25pt,
    arc=3pt,
    left=16pt,right=16pt,top=14pt,bottom=14pt,
    borderline north={3.2pt}{0pt}{StageBlue},
    borderline west={4pt}{0pt}{StageBlue},
    borderline south={1.25pt}{0pt}{StageBlue},
    drop shadow={black!25!white},
  },
  charterRule/.style={
    enhanced,
    colback=LightGray,
    colframe=RuleGray,
    boxrule=0.85pt,
    arc=3pt,
    left=12pt,right=12pt,top=10pt,bottom=10pt,
    borderline west={3pt}{0pt}{StageBlue},
  },
  charterBadge/.style={
    enhanced,
    colback=StageBlue,
    coltext=white,
    colframe=StageBlueDark,
    boxrule=0.6pt,
    arc=2pt,
    left=6pt,right=6pt,top=3pt,bottom=3pt,
  }
}
\newtcbox{\Badge}{nobeforeafter,charterBadge}

% ---------- Lists ----------
\setlist[itemize]{leftmargin=1.5em, itemsep=0.25em, topsep=0.25em}
\setlist[enumerate]{leftmargin=1.8em, itemsep=0.25em, topsep=0.25em}

% ---------- TOC ----------
\usepackage{tocloft}
\setcounter{tocdepth}{3}
\setlength{\cftbeforesecskip}{3pt}
\setlength{\cftbeforesubsecskip}{2pt}
\setlength{\cftbeforesubsubsecskip}{0pt}
\renewcommand{\cftsecfont}{\bfseries}
\renewcommand{\cftsecpagefont}{\bfseries}
\renewcommand{\cftsubsecfont}{\normalfont}
\renewcommand{\cftsubsecpagefont}{\normalfont}
\renewcommand{\cftsubsubsecfont}{\normalfont}
\renewcommand{\cftsubsubsecpagefont}{\normalfont}
\setlength{\cftsecindent}{0pt}
\setlength{\cftsubsecindent}{1.25em}
\setlength{\cftsubsubsecindent}{2.8em}
\setlength{\cftsecnumwidth}{2.1em}
\setlength{\cftsubsecnumwidth}{2.8em}
\setlength{\cftsubsubsecnumwidth}{3.6em}

% Remove the built-in TOC heading so only the custom heading is shown
\makeatletter
\renewcommand{\tableofcontents}{\@starttoc{toc}}
\makeatother

% ---------- Header/Footer: page number only ----------
\usepackage{fancyhdr}
\pagestyle{fancy}
\fancyhf{}
\renewcommand{\headrulewidth}{0pt}
\renewcommand{\footrulewidth}{0pt}
\fancyfoot[C]{\thepage}

% ---------- Section formatting ----------
\usepackage{titlesec}

\titleformat{\section}
  {\Large\bfseries\color{Ink}}
  {\thesection}{0.6em}{}
  [\vspace{0.25em}\textcolor{StageBlue}{\rule{\linewidth}{0.8pt}}]

\titleformat{\subsection}
  {\large\bfseries\color{Ink}}
  {\thesubsection}{0.6em}{}

\titleformat{\subsubsection}
  {\normalsize\bfseries\color{Ink}}
  {\thesubsubsection}{0.6em}{}

\titlespacing*{\section}{0pt}{0.95em}{0.35em}
\titlespacing*{\subsection}{0pt}{0.65em}{0.25em}
\titlespacing*{\subsubsection}{0pt}{0.55em}{0.2em}

% ---------- Table engine ----------
\usepackage{tabularray}
\UseTblrLibrary{booktabs}

% ---------- tabularray: GLOBAL table treatment ----------
\SetTblrInner{
  cells = {fg=black, bg=white, valign=t},
}

% ---------- Header helpers ----------
\newcommand{\Hdr}[1]{\shortstack[c]{\strut #1\strut}}

% ---------- Lines ----------
\newcommand{\TitleLine}{\textcolor{StageBlue}{\rule{0.98\linewidth}{0.95pt}}}
\newcommand{\SubTitleLine}{\textcolor{RuleGray}{\rule{0.98\linewidth}{0.65pt}}}

% ---------- Continued on next page ----------
\newcommand{\ContinuedOnNextPageInline}{%
  \par
  \vfill
  \noindent\textcolor{RuleGray}{\rule{\linewidth}{1pt}}\vspace{-2pt}
  \noindent\hfill\textit{Continued on next page}\par\vspace{-10pt}
  \noindent\textcolor{RuleGray}{\rule{\linewidth}{1pt}}
  \clearpage
}

% ---------- PDF hyperlinks ----------
\usepackage[hidelinks]{hyperref}
\hypersetup{
  colorlinks=false,
  pdfborder={0 0 0},
  linktoc=all,
  pdftitle={\DocTitle},
  pdfauthor={\ProgramName}
}

% =========================================================
\begin{document}

\section*{7.20 — Personal Hygiene, Health, and Sanitation Items}
\phantomsection
\addcontentsline{toc}{section}{7.20 — Personal Hygiene, Health, and Sanitation Items}

\subsection*{7.20.1 Purpose \& Scope}
\phantomsection
\addcontentsline{toc}{subsection}{7.20.1 Purpose \& Scope}

7.20 covers consumables and small tools that preserve hygiene, sanitation, and basic health maintenance for training continuity. It targets preventable failure points: chafing, skin irritation, persistent odor, fungal conditions, small wounds worsening, and wet-gear loops that drive missed sessions.

\subsection*{7.20.2 Governance and Safety Boundaries}
\phantomsection
\addcontentsline{toc}{subsection}{7.20.2 Governance and Safety Boundaries}

\begin{itemize}
\item Hygiene is an adherence and durability control; it prevents preventable issues but does not replace medical care.
\item Non-medical escalation triggers (seek evaluation posture):
  \begin{itemize}
  \item infection signs (spreading redness/heat, swelling, increasing pain, drainage),
  \item fever/systemic illness signs,
  \item worsening wounds or non-healing lesions,
  \item numbness/tingling or rapidly worsening symptoms.
  \end{itemize}
\item No instruction for self-treating serious conditions; keep guidance preventive and operational.
\item Blister boundary:
  \begin{itemize}
  \item 7.11 owns the blister kit and blister protocol items.
  \item 7.20 may provide hygiene-adjacent barrier consumables and cleaning posture that support foot integrity but must not duplicate full blister kit contents.
  \end{itemize}
\item Nail care ownership:
  \begin{itemize}
  \item clippers/file are owned by 7.20 and are cross-referenced by 7.11.
  \end{itemize}
\end{itemize}

\subsection*{7.20.3 Tier Doctrine — Applied to Hygiene/Sanitation}
\phantomsection
\addcontentsline{toc}{subsection}{7.20.3 Tier Doctrine — Applied to Hygiene/Sanitation}

\begin{itemize}
\item Tier 0: household-only substitutes and baseline hygiene posture; minimal feasibility controls.
\item Tier 1: minimum deployable hygiene kit to support six sessions/week continuity, skin integrity, and gear cleanliness.
\item Tier 2: expanded travel continuity and higher-volume support (more redundancy, better portability, stronger skin integrity support).
\item Tier 3: overbuilt redundancy and staging for high-consequence phases and travel.
\end{itemize}

\subsection*{7.20.4 Selection Criteria \& Fit Checks}
\phantomsection
\addcontentsline{toc}{subsection}{7.20.4 Selection Criteria \& Fit Checks}

\begin{itemize}
\item Skin sensitivity:
  \begin{itemize}
  \item choose gentle posture first; discontinue and switch if rash/irritation occurs.
  \end{itemize}
\item Portability:
  \begin{itemize}
  \item travel bottles must be leakproof; kits must fit go-bag posture (7.10 staging).
  \end{itemize}
\item Clothing and foot interface linkage:
  \begin{itemize}
  \item chafe risk often originates from fit/fabric (7.15) and sock/footwear interface (7.11); address those first.
  \end{itemize}
\item Fit checks:
  \begin{itemize}
  \item bandages adhere without tearing skin,
  \item anti-chafe product does not irritate and does not cause slipping hazards.
  \end{itemize}
\end{itemize}

\subsection*{7.20.5 Setup, Safety, and Anchoring}
\phantomsection
\addcontentsline{toc}{subsection}{7.20.5 Setup, Safety, and Anchoring}

\begin{itemize}
\item Routine posture:
  \begin{itemize}
  \item post-session wash and dry; avoid staying in wet clothes.
  \item launder clothing and socks on a cadence that prevents repeated damp reuse.
  \end{itemize}
\item Chafe posture:
  \begin{itemize}
  \item adjust fit/fabric first; apply anti-chafe consumables only to known hotspots and stop if irritation occurs.
  \end{itemize}
\item Minor wound posture:
  \begin{itemize}
  \item cover and monitor; keep clean; seek evaluation if worsening signs occur.
  \end{itemize}
\item Nail care posture:
  \begin{itemize}
  \item maintain nails trimmed and filed; prevent jagged edges that tear skin.
  \end{itemize}
\item Staging posture:
  \begin{itemize}
  \item keep hygiene kit staged with training kit (cross-reference 7.10).
  \end{itemize}
\end{itemize}

\subsection*{7.20.6 Maintenance, Replacement, and Storage}
\phantomsection
\addcontentsline{toc}{subsection}{7.20.6 Maintenance, Replacement, and Storage}

\begin{itemize}
\item Restocking cadence:
  \begin{itemize}
  \item maintain minimum inventory to avoid “no supplies” failures.
  \end{itemize}
\item Replacement cues:
  \begin{itemize}
  \item expired supplies, degraded adhesives, cracked travel bottles, persistent contamination/odor that cannot be resolved.
  \end{itemize}
\item Storage:
  \begin{itemize}
  \item dry and sealed; avoid moisture exposure; keep kits accessible and staged.
  \end{itemize}
\end{itemize}

\subsection*{7.20.7 Substitution Rules — Minimal-Path Alternatives}
\phantomsection
\addcontentsline{toc}{subsection}{7.20.7 Substitution Rules — Minimal-Path Alternatives}

\begin{itemize}
\item If hygiene kit is incomplete:
  \begin{itemize}
  \item use soap + water baseline and basic laundry access posture as minimum substitutes.
  \item prioritize cleanliness and drying over any specialty products.
  \end{itemize}
\item If anti-chafe product is not tolerated:
  \begin{itemize}
  \item remove it, adjust clothing fit/fabric (7.15), adjust sock interface (7.11), and downshift exposure if needed.
  \end{itemize}
\item If minor wounds worsen:
  \begin{itemize}
  \item stop and seek evaluation posture; do not self-treat beyond cover/monitor.
  \end{itemize}
\item Missing items:
  \begin{itemize}
  \item do not improvise with unsafe substitutes; stage replacements promptly.
  \end{itemize}
\end{itemize}

\subsection*{7.20.8 Phase Integration Map — Required}
\phantomsection
\addcontentsline{toc}{subsection}{7.20.8 Phase Integration Map — Required}

\begingroup
\scriptsize
\setlength{\tabcolsep}{2pt}
\renewcommand{\arraystretch}{1.10}

\begin{longtblr}{
  width=\linewidth,
  rowhead=1,
  colspec = {
    Q[l,wd=1.05in]
    Q[l,wd=1.55in]
    X[l,2.75]
    X[l,3.65]
  },
  hlines,
  vlines,
  cells = {hsep=2pt, vsep=6pt, halign=l, valign=t},
  row{1} = {bg=StageBlue!15, fg=black, font=\bfseries, halign=l, valign=m},
  row{odd} = {bg=white},
  row{even} = {bg=LightGray},
}
\Hdr{Phase} &
\Hdr{Required Tier (from Stage 6)} &
\Hdr{What Changes / Becomes Mandatory} &
\Hdr{Notes} \\
Phase 00 &
T1 (E) &
Tier 1 hygiene system is required by Phase 00 (E): body wash, shampoo, hand soap, deodorant, oral-care basics, nail care, towels, laundry posture, coverage supplies, and anti-chafe posture &
Phase 00 (E): start-from-zero hygiene prevents early skin breakdown and illness-driven drift. \\
Phase 01 &
T1 (End) &
Maintain Tier 1; required restock discipline increases as volume rises &
Phase 01 (End): maintain staged supplies; prevent “no clean socks/no supplies” misses. \\
Phase 02 &
T1 (End) &
Maintain Tier 1; required skin integrity posture supports higher workload &
Phase 02 (End): treat persistent irritation as a downshift trigger; seek evaluation if worsening. \\
Phase 03 &
T1 (End) &
Maintain Tier 1; required attention to foot moisture and chafe increases as ruck begins &
Phase 03 (End): cross-reference 7.11 blister protocol; 7.20 supports barrier/cleaning posture only. \\
Phase 04 &
T1 (End) &
Maintain Tier 1; required wet-gear and swim-related hygiene posture increases &
Phase 04 (End): aquatics introduces wet gear; drying/logistics via 7.10; consumables here support hygiene. \\
Phase 05 &
T1 (End) &
Maintain Tier 1; required prevention of damp clothing loops increases with endurance volume &
Phase 05 (End): laundry cadence and dry storage prevent skin irritation and odor loops. \\
Phase 06 &
T1 (End) &
Maintain Tier 1; Tier 2 items become optional continuity aids &
Phase 06 (End): travel kit posture is higher ROI but Stage 6 does not elevate 7.20 beyond Tier 1. \\
Phase 07 &
T1 (End) &
Maintain Tier 1; higher density increases restock pressure &
Phase 07 (End): staged redundancy may be useful but remains optional unless Stage 6 changes. \\
Phase 08 &
T1 (End) &
Maintain Tier 1; load tolerance increases skin/foot integrity consequence &
Phase 08 (End): prevent small issues from becoming gate-disrupting events. \\
Phase 09 &
T1 (End) &
Maintain Tier 1; high \textbf{CAL} volume increases sweat and friction risk &
Phase 09 (End): ensure anti-chafe education-only posture and laundry posture are stable. \\
Phase 10 &
T1 (End) &
Maintain Tier 1; multi-domain weeks increase consequence of illness/skin breakdown &
Phase 10 (End): enforce no novelty and stable routines. \\
Phase 11 &
T1 (End) &
Maintain Tier 1; recovery sensitivity increases consequence of hygiene failures &
Phase 11 (End): treat repeated irritation as a risk flag requiring downshift. \\
Phase 12 &
T1 (End) &
Maintain Tier 1; Tier 3 redundancy remains optional &
Phase 12 (End): if redundancy adopted, stage via 7.10; still not required unless Stage 6 changes. \\
Phase 13 &
T1 (End) &
Maintain Tier 1; consolidation demands low variance and stable hygiene routines &
Phase 13 (End): prevent last-minute disruptions from preventable hygiene failures. \\
\end{longtblr}

\endgroup

7.20.8 Phase Integration Map Notes:

\begin{itemize}
\item If any phase artifact requires 7.20 Tier 2/3 as mandatory earlier than Stage 6 timing, requires Stage 8 review.
\item If any artifact duplicates the 7.11 blister kit within 7.20, requires Stage 8 review.
\end{itemize}

\subsection*{7.20.9 Risks, Contraindications, and Red Flags}
\phantomsection
\addcontentsline{toc}{subsection}{7.20.9 Risks, Contraindications, and Red Flags}

\begin{itemize}
\item Risks:
  \begin{itemize}
  \item skin irritation and breakdown from damp clothing loops and poor laundry cadence,
  \item chafing leading to altered mechanics,
  \item minor wounds worsening due to poor cleaning and coverage,
  \item fungal issues from prolonged damp conditions.
  \end{itemize}
\item Red flags requiring stop and seek evaluation posture (non-medical):
  \begin{itemize}
  \item spreading redness/heat, increasing swelling, drainage, fever/systemic illness signs,
  \item rapidly worsening pain or gait-altering pain associated with skin lesions.
  \end{itemize}
\item Contraindication posture:
  \begin{itemize}
  \item discontinue any product causing rash/irritation; switch to gentler posture; seek evaluation if persistent.
  \end{itemize}
\item Requires Stage 8 review triggers:
  \begin{itemize}
  \item any timing mismatch requiring 7.20 earlier than Stage 6,
  \item any content drifting into medical wound-care protocols or treatment claims.
  \end{itemize}
\end{itemize}

\end{document}